\section{Architectural Approach and System Overview}
\label{sec:design-style}
% Authoring brief for this section lives in section.md (§5.2).

The drivers of \cref{sec:design-drivers} specify the forces the architecture must reconcile, but not the form it should take.
This section fixes that form in two steps, moving from the system in outline to the style that governs its internal structure, before any single concern is opened in \cref{sec:design-modules}.
We first present the platform at two levels of abstraction, the system in its environment and the runtime parts beneath it, and then name and justify the architectural style that the dominant driver D1 demands.

At the highest level of abstraction the platform is a single, researcher-operated system that mediates a participant's reading session, shown in its environment in \cref{fig:design-context}.
The platform sits between the two people it serves: the participant, who reads, and the researcher, who operates and monitors the session through a separate console (\cref{sec:stakeholders}).
Two further elements meet the platform at its boundary rather than forming part of its core: the eye tracker (a Tobii device in this work) is the sensing interface that supplies the gaze signal, and an external decision provider may attach through a well-defined API to supply the adaptive intelligence that this thesis treats as out of scope (C3, \cref{sec:constraints}).
We frame this outermost view as a system context in the sense of the C4 model \cite{c4model}; the runtime parts beneath it and the internal module structure are opened in the sections that follow.
This outermost view fixes the system boundary, and with it the division of responsibility that the rest of the chapter elaborates: the platform owns the capture of gaze, its analysis, the execution of interventions, and the session record, whereas it deliberately owns neither the internals of the eye-tracking hardware nor the decision intelligence behind the provider seam.

\begin{figure}[htbp]
\centering
\includesvg[width=\linewidth,inkscapeopt={--export-text-to-path}]{Chapters/05_SystemDesign/architecture_overview}
\caption{Overview of the Reading the Reader platform as a researcher--participant adaptive loop. The participant reads on one screen while a Tobii eye tracker at the foot of that screen senses gaze; the backend turns the gaze stream into oculomotor events, adapts the participant's text, and mirrors the gaze and derived analysis to the researcher's console, where the session is steered and proposed interventions are approved or rejected. An external decision provider may attach through a well-defined API, to which the platform sends gaze and context and from which it receives intervention commands (out of scope, C3); every session is written to a reproducible record.}
\label{fig:design-context}
\end{figure}

Beneath this overview, the platform resolves into a small set of runtime parts whose relationships the rest of the chapter elaborates.
The researcher's console and the participant's reading surface are two distinct surfaces of a single web front-end rather than a shared screen, so that the participant sees only the text while the researcher retains a mirror of that view together with the controls around it (D6).
Neither surface communicates with the other: both talk only to the backend (the application core), which carries the core logic and mediates between the two surfaces, the eye tracker, and the store.
Two channels connect the front-end to the backend: a request and response channel carries discrete operator commands and setup data, whereas a separate high-frequency channel carries the live gaze stream, a separation we introduce here and justify against the responsiveness budget D2 in \cref{sec:design-loop}.
Session state is written to a local, file-backed store rather than an external database, a choice we introduce here and defend on grounds of reproducibility and scope in \cref{sec:design-data,sec:design-decisions}.
The external decision provider attaches through a single, well-defined API rather than being embedded in the backend, which is the structural expression of the out-of-scope constraint C3.
The concrete technologies (a Next.js and React front-end, an ASP.NET Core backend) are justified, together with the alternatives considered, in \cref{sec:technology-selection}.

The two front-end surfaces differ in what they present.
The participant's reading surface shows only the text, whereas the researcher's console mirrors that text and overlays the live gaze together with the fixation heat map and saccades that the core derives, so that the researcher sees not merely where the participant looks but how they read.
The same console carries the controls that steer the session, the choice between advisory and autonomous operation and the approval of individual interventions, while an external decision provider may attach through the out-of-scope seam.
This pairing of a clean reading surface for the participant with an analytic, controllable console for the researcher is the platform's central user-facing structure, whose behaviour in motion and whose control by the researcher are taken up in \cref{sec:design-loop,sec:design-extensibility}.

The internal structure of the backend follows one organising principle, namely that dependencies point inward, from volatile detail towards stable abstraction.
Hardware access, transport, persistence, and presentation are treated as replaceable detail that depends on the core through abstractions, while the core itself depends on nothing outward \cite{martin2017cleanarch, martin1996dip}.
Expressed in the vocabulary of ports and adapters, each external concern meets the core at a port, an interface owned by the core, and each concrete technology is an adapter that the infrastructure supplies behind that port \cite{cockburn2005hexagonal}.
This is one idea seen from two angles, the dependency rule describing the direction and the ports-and-adapters metaphor describing the seam, and we use both terms in the sections that follow.

This style is not adopted for its own sake, but because it is the most direct structural response to the dominant driver, modular separability D1.
Making every boundary an interface, and forbidding the core from depending on any concrete technology, is precisely what allows sensing, analysis, decision, and intervention to be replaced independently, which is the claim that RQ1 puts to the test.
The alternative is the shape taken by the earlier Reading the Reader prototypes, in which sensing, adaptation, and presentation were realised as a single tightly coupled application \cite{refsgaard2024rtr, pereira2024typography, palinec2024reactive}; that coupling is the limitation identified in \cref{sec:gaps-opportunities} and the reason a new architecture is warranted.
We defer the full comparison of this and the other structural options to the decision register of \cref{sec:design-decisions}, and turn first to the concerns upon which the inward-pointing dependencies converge.
